\section{Ordinary cooperations and finite tests}\label{sec:cooperations}

The upper exponent bound requires the cohomological line of
Section~\ref{sec:adic} after every finite ordinary test.
We therefore identify the tested Amitsur coefficient complex, including
its actual coface and codegeneracy maps.
Fix a positive height $a$, and write
\[
 \begin{gathered}
 E=E_a,\qquad k_a=\F_{p^a},\qquad
 O=O_a=W(k_a)[[u_1,\ldots,u_{a-1}]],\\
 O_*=(E_a)_*=O[u^{\pm1}],\quad |u|=-2,\qquad
 \m=\m_a=(p,u_1,\ldots,u_{a-1}).
 \end{gathered}
\]
Use the coherent action of $\G_a=\St_a\rtimes\Gal(k_a/\F_p)$ by
commutative ring automorphisms, whose special-fiber action is the
prescribed Honda stabilizer action
\cite[Section~7, Proposition~7.3, Equation~(7.3), and
Corollaries~7.6--7.7]{gh2004}.
The local category $\mathcal C_a=\mathrm{Sp}_{K(a)}$ has unit
$B_a=L_{K(a)}S$ and tensor
$U\otimes_a V=L_{K(a)}(U\wedge V)$.
All unadorned smash products and function spectra below are ordinary.

\begin{proposition}[The finite-test coefficient complex]\label{wup:cochains}
For \(r\geq0\), let
\[
 Z_r=L_{K(a)}(E^{\wedge(r+1)}).
\]
For every finite ordinary \(p\)-local spectrum \(W\), there are
natural graded isomorphisms
\begin{equation}\label{wup:tested-functions}
 \pi_*(Z_r\wedge W)
 \cong C_{\cts}(\G_a^r,\pi_*(E\wedge W)).
\end{equation}
The first \(E\)-factor acts by constant coefficients.
The two-factor evaluation is multiplication after \(1\wedge g\).
In the cumulative coordinates \(1,g_1,g_1g_2,\ldots,g_1\cdots g_r\)
on the factors, these isomorphisms identify every coface and
codegeneracy of the unit Amitsur object with the continuous
inhomogeneous group-cochain operator.

Consequently, apply the exact ordinary functor \(F(V_0,-)\), for a
finite ordinary \(p\)-local spectrum \(V_0\), to the full unit
Amitsur tower for \(E_a\) in \(\mathcal C_a\). In the page convention
of Proposition~\ref{prop:pages}, its second page is
\begin{equation}\label{wup:cochain-page}
 E_2^{s,q}
 =H^s_{\cts}\bigl(\G_a,\pi_q(E_a\wedge D V_0)\bigr),
 \qquad D V_0=F(V_0,S).
\end{equation}
Here \(H^s_{\cts}\) means the cohomology of continuous cochains with
the adic topology on the indicated coefficient module.
\end{proposition}

We prove the comparison by first forming ordinary cooperations, then
completing their flat coefficient modules, and finally applying the
finite test. The next three lemmas justify these steps. Mixed
cooperations and their completion will also be used in
Section~\ref{sec:coherent-trace}.
The exact-homology input originates in Landweber's work
\cite{landweber1976}; we use the precise cooperation formulations cited
below, with $BP$ the $p$-local Brown--Peterson spectrum.

\begin{lemma}[Ordinary mixed cooperations]\label{wup:mixed-flat}
For positive heights \(a,b\), the ordinary smash product has even
homotopy
\[
 \pi_*(E_a\wedge E_b)\cong
 (E_a)_*\otimes_{BP_*}BP_*BP\otimes_{BP_*}(E_b)_*.
\]
It is flat over \((E_a)_*\) through the first factor.
\end{lemma}

\begin{proof}
Apply the ordinary Landweber cooperation calculation to the $p$-typical
flat Hopf algebroid $(BP_*,BP_*BP)$. This gives the displayed $BP_*$
formula. It is the $p$-typical form of the $MU_*$ formulas displayed in
\cite[Section~15, Proposition~15.1 and Corollary~15.4]{rezk1998} and
\cite[Lemma~4.4]{hovey2004}; Morava $E$-theory over $BP$ is included in
\cite[Example~0.1(d)]{hs-comodules2003}.
The two tensor products use the two different unit maps into \(BP_*BP\).
For a Landweber-exact \(BP_*\)-algebra \(C\), the module
\(BP_*BP\otimes_{BP_*}C\) is flat over the other \(BP_*\)-action.
Indeed, tensoring with it factors as the exact cofree-comodule functor
followed by the exact Landweber functor on comodules. This is also the
opposite-unit criterion of
\cite[Lemma~2.2]{hs-comodules2003}, with conjugation exchanging the
units when necessary. Base change along \(BP_*\to(E_a)_*\) proves the
asserted flatness. The displayed coefficient and cooperation rings are
even graded, which proves evenness. In particular this argument does
not require \((E_b)_*\) to be flat over \(BP_*\).
\end{proof}

The ring \(O\) is complete, Noetherian, regular, and local, with finite
residue field. If \(W\) is a finite ordinary \(p\)-local spectrum, then
\[
 M_*=\pi_*(E\wedge W)
\]
is finitely generated over \(O_*\). This follows by finite cofiber
constructions and retracts from the assertion for the sphere.
Periodicity identifies graded module questions with the two parities
over \(O\). Regularity therefore supplies a finite resolution
\cite[Tag~00O7]{stacks} of
\(M_*\) by finite graded free \(O_*\)-modules. Neither parity of
\(M_*\) is required to vanish.

\begin{lemma}[Base change and continuous functions]\label{wup:base-change}
Let \(D_0\) be an ordinary right \(E\)-module for which
\(\pi_*D_0\otimes_{O_*}(-)\) is exact on finitely generated graded
\(O_*\)-modules. The natural comparison
\begin{equation}\label{wup:relative-base-change}
 \pi_*D_0\otimes_{O_*}\pi_*(E\wedge W)
 \longrightarrow \pi_*(D_0\wedge W)
\end{equation}
is an isomorphism for every ordinary \(p\)-local spectrum \(W\).

For every profinite set \(T\), the functor \(C_{\cts}(T,-)\) is exact
on finitely generated \(O\)-modules with their \(\m\)-adic topologies.
For every finitely generated graded \(O_*\)-module \(M_*\), the
natural map is an isomorphism
\begin{equation}\label{wup:function-tensor}
 C_{\cts}(T,O_*)\otimes_{O_*}M_*
 \xrightarrow{\ \cong\ } C_{\cts}(T,M_*).
\end{equation}
The functions and the displayed graded isomorphism are understood
degreewise.
\end{lemma}

\begin{proof}
First suppose \(W\) finite. The equality
\[
 D_0\wedge W\simeq D_0\wedge_E(E\wedge W)
\]
is an equality of ordinary derived constructions. The comparison is
the balanced homotopy cross product for this relative smash, so it is
natural in both variables before choosing any resolution.
Realize a finite graded free surjection onto \(\pi_*(E\wedge W)\) by an \(E\)-module
map from a finite wedge of suspensions of \(E\). Surjectivity makes the
homotopy of its fiber the first syzygy. Repeat along a finite free
resolution. At the last stage, lifts of a free basis give an
equivalence from a finite free \(E\)-module, since they give an
isomorphism on all homotopy groups. Relative smash with \(D_0\) is
exact. Induction back through these fiber sequences, using the
assumed tensor exactness on each finitely generated syzygy, proves
\eqref{wup:relative-base-change} for finite \(W\).

Every ordinary \(p\)-local spectrum is a filtered homotopy colimit of
finite ordinary \(p\)-local spectra. Both sides of
\eqref{wup:relative-base-change} preserve those colimits: ordinary
smash and homotopy groups do so, as does tensoring a module with
\(\pi_*D_0\). Naturality extends the comparison to arbitrary \(W\).

For the assertion about functions, let \(M\to N\) be a surjection of
finitely generated \(O\)-modules. For every \(j\geq1\) the map
\[
 M/\m^{j+1}M\longrightarrow
 (M/\m^jM)\mathbin{\times}_{N/\m^jN}(N/\m^{j+1}N)
\]
is surjective. To see this, lift the first component to \(M\); its
error against the second component lies in \(\m^jN\), and this error
lifts from \(\m^jM\). Each quotient is a finite set. Given a continuous
map \(T\to N\), first lift its reduction modulo \(\m\), and then use
these surjections to lift successively while preserving the previous
choice. A lift on each finite value set is continuous. The compatible
system defines a continuous lift to \(M\), since finitely generated
modules are complete. Kernels carry their subspace topologies:
the relevant modules are compact Hausdorff, and the induced
continuous bijection from a kernel to its image is a homeomorphism.
It follows that \(C_{\cts}(T,-)\) is exact on these modules. Apply this
exact functor to a finite presentation of \(M\), and compare with
tensoring the same presentation by \(C_{\cts}(T,O)\). The comparison
is an isomorphism on the finite free terms and hence on \(M\).
Applying this argument in each parity proves
\eqref{wup:function-tensor}. In particular
\(C_{\cts}(T,O_*)\otimes_{O_*}(-)\) has precisely the tensor exactness
used in the first assertion.
\end{proof}

\begin{lemma}[Flat completion and reduction]\label{wup:completion-reduction}
Let \(N\) be an ordinary \(E\)-module with flat graded homotopy
\(U_*=\pi_*N\). Its actual \(K(a)\)-localization map induces
\[
 U_*\longrightarrow\pi_*L_{K(a)}N
 \cong\widehat{U_*}:=\varprojlim_j U_*/\m^jU_*.
\]
For a flat \(O\)-module \(U\) and every \(j\geq1\),
\[
 \widehat U/\m^j\widehat U\cong U/\m^jU.
\]
Both assertions apply degreewise to the periodic graded modules above.
\end{lemma}

\begin{proof}
The first assertion is the localization theorem for ordinary Morava
\(E\)-modules: \(K(a)\)-localization is derived \(\m\)-completion,
and flat homotopy has no higher derived completion. The natural
localization map induces the ordinary completion map
\cite[Propositions~3.4, 3.6 and Corollary~3.8]{rezk2009}.

Here is the algebraic reduction argument, including its needed
inverse-limit step. Choose a resolution \(P_\bullet\to N_0\) of a
finitely generated \(O\)-module \(N_0\) by finite free modules.
The tower \(P_\bullet\otimes_O U/\m^lU\) is degreewise surjective.
By flatness of \(U\), its first homology is
\[
 \operatorname{Tor}_1^O(N_0,O/\m^l)\otimes_O U.
\]
This homology tower is pro-zero. Indeed, for a finite first
presentation \(0\to K\to P_0\to N_0\to0\), the corresponding Tor group
is \((K\cap\m^lP_0)/\m^lK\). Artin--Rees
\cite[Tag~00IN]{stacks} gives a \(d\) such that
\(K\cap\m^lP_0\subseteq\m^{l-d}K\) for \(l\geq d\).
For any fixed \(j\), a transition from \(l\geq j+d\) to \(j\)
is therefore zero. Tensoring these zero transitions with \(U\)
preserves this conclusion.

The homology exact sequence for the countable inverse limit of this
degreewise-surjective tower now gives
\[
 N_0\otimes_O\widehat U
 \cong\varprojlim_l(N_0\otimes_O U/\m^lU);
\]
the possible \(\varprojlim^1\) of the first homology is zero because
that tower is pro-zero. Finite freeness of the \(P_i\) identifies
their tensor with \(\widehat U\) with their inverse-limit tensor.
Take \(N_0=O/\m^j\); for \(l\geq j\) the modules on the right are
\(U/\m^jU\). This proves the reduction formula.
\end{proof}

\begin{proof}[Proof of Proposition~\ref{wup:cochains}]
The completed cooperation theorem gives
\begin{equation}\label{wup:completed-cooperation}
 \pi_*L_{K(a)}(E\wedge E)
 \cong C_{\cts}(\G_a,O_*).
\end{equation}
At \(g\), the map is \(\mu(1\wedge g)\); the first factor acts by
constant scalars. This is the evaluation map of
\cite[Section~1 and Theorems~4.11 and 6.8]{hovey2004}.
The equivalent convention in
\cite[Proposition~2.2, Equation~(2.3), and Equations~(2.5)--(2.7)]{dh2004}
uses \(\mu(g^{-1}\wedge1)\); exchanging the factors and inverting
the group coordinate gives the convention above. The coefficient
topology is the adic one
\cite[Remark~1.3]{dh2004}.

Put \(P_*=\pi_*(E\wedge E)\), regarded as an \(O_*\)-module through
the first factor. Lemma~\ref{wup:mixed-flat} makes \(P_*\) flat.
The actual localization map, the reduction in
Lemma~\ref{wup:completion-reduction}, and
\eqref{wup:function-tensor} applied to \(O_*/\m^jO_*\) give
\begin{equation}\label{wup:ordinary-reductions}
 P_*/\m^jP_*
 \cong C_{\cts}(\G_a,O_*/\m^jO_*).
\end{equation}
This isomorphism is induced by the actual ordinary evaluation
\(E\wedge E\xrightarrow{1\wedge g}E\wedge E\xrightarrow{\mu}E\)
followed by reduction. It is not just an abstract isomorphism of
the two reduction modules.

To form higher powers, use a common first factor in the ordinary
relative smash products
\[
 (E_0\wedge E_1)\wedge_{E_0}\cdots
       \wedge_{E_0}(E_0\wedge E_r)
 \simeq E_0\wedge E_1\wedge\cdots\wedge E_r.
\]
Each \(E_i\) denotes the same spectrum \(E\); the subscripts mark
positions. Iterating \eqref{wup:relative-base-change}, including
its extension to arbitrary ordinary \(W\), identifies the homotopy
of this ordinary smash with \(P_*^{\otimes_{O_*}r}\).
The comparison inserts each pair into the first and one other
position and multiplies at the common position. These modules are
flat, since \(P_*\) is flat.

In each degree the coefficient ring \(O_*/\m^jO_*\) is finite.
Continuous functions to it are unions of function modules on finite
quotient sets of \(\G_a\). Tensors of the function modules on finite
sets are the function modules on their products, and every
continuous function on \(\G_a^r\) to a finite set factors through
finite quotients in the coordinates. It follows from
\eqref{wup:ordinary-reductions} that the reductions of
\(P_*^{\otimes_{O_*}r}\) are
\(C_{\cts}(\G_a^r,O_*/\m^jO_*)\). Apply
Lemma~\ref{wup:completion-reduction} only now, to this ordinary
\((r+1)\)-fold smash. The result is
\begin{equation}\label{wup:power-functions}
 \pi_*Z_r\cong C_{\cts}(\G_a^r,O_*).
\end{equation}
For \(r=0\) this is the locality of \(E\).

The maps in \eqref{wup:power-functions} remain the evaluations on
the non-first factors followed by multiplication. In fact each
ordinary evaluation extends through localization because its target
\(E\) is local. On homotopy its first-factor \(O_*\)-linearity makes
it adically continuous. Lemma~\ref{wup:completion-reduction}
makes the ordinary image dense, and the reductions above show that
the extension agrees with the displayed continuous function.
Thus the construction preserves the evaluation maps as well as the
coefficient modules.

For finite \(W\), the ordinary \(E\)-module \(Z_r\) satisfies the
exactness assumption in Lemma~\ref{wup:base-change}, by
\eqref{wup:power-functions} and \eqref{wup:function-tensor}.
Those comparisons give \eqref{wup:tested-functions}. This step
allows any finitely generated \(\pi_*(E\wedge W)\), including
torsion and both parities; no flatness of that module has been used.

The terms of the local unit Amitsur object are exactly \(Z_r\).
For finite \(V_0\), ordinary duality identifies
\(F(V_0,Z_r)\simeq Z_r\wedge D V_0\), so the previous formula
computes each cosimplicial homotopy group. Relabel the independent
evaluations by their cumulative coordinates
\[
 1,\quad g_1,\quad g_1g_2,\quad\ldots,\quad g_1\cdots g_r.
\]
The cochain identities below can be checked on ordinary evaluations before
localization. For finite $W$, their target $E\wedge W$ is $K(a)$-local,
since finite dualizable smash preserves local objects. They therefore
extend through localization and remain identities on the tested
coefficients. Insertion of the first unit acts on the coefficient; an
interior unit combines two consecutive coordinates, and
insertion of the last unit drops the last coordinate. Explicitly,
for a degree-\(r\) cochain \(c\),
\[
\begin{aligned}
 (\delta^0c)(g_1,\ldots,g_{r+1})
   &=g_1c(g_2,\ldots,g_{r+1}),\\
 (\delta^ic)(g_1,\ldots,g_{r+1})
   &=c(g_1,\ldots,g_i g_{i+1},\ldots,g_{r+1})
       &&(1\leq i\leq r),\\
 (\delta^{r+1}c)(g_1,\ldots,g_{r+1})
   &=c(g_1,\ldots,g_r).
\end{aligned}
\]
Multiplication of adjacent factors gives each codegeneracy, which
inserts the identity in the corresponding coordinate. These are all
the inhomogeneous cochain operators, including the degeneracies of
the full cosimplicial object. The first unit insertion also displays
the orbit function of every coefficient. Those functions are
continuous, and reduction gives continuous actions on the finite
invariant quotients by powers of \(\m_a\). The \(E_2\)-page of the
tower of finite partial totalizations is cosimplicial cohomology
of these homotopy groups, which is \eqref{wup:cochain-page}.
\end{proof}
